\section*{Materials and methods}

\subsection*{Study overview}
Fig~\ref{fig:overview} shows how the parts of the study fit together. An
agent ran eight interpretability pipelines on MaxToki-217M by following
the specifications, over several working sessions. Then, in a single
session on the last day, it audited its outputs with the checklist (the
\emph{deployment} and the \emph{in-session audit}). Everything after that was done separately and is the evaluation:
the independent verification that built the error ledger, the three
controlled studies, and the corrected MaxToki-217M analyses. The
controlled studies used fresh agents that had no access to the ledger,
the answer keys, this manuscript or the project's working notes.

\begin{figure}[H]
\caption{\textbf{Study overview.} Left: what was deployed. Each
interpretability method is written as a markdown specification; an agent
follows the specification to run the analysis on MaxToki-217M over
several working sessions and writes run summaries. Then, in one session
on the last day, the agent applies the ten-pattern checklist to the
outputs of all eight pipelines and repairs what it finds. Right: how the deployment was
evaluated. Independent verification (number tracing by code,
re-derivation of statistics from stored per-record outputs, unit tests of
intervention code, an input-encoding audit) and blind re-audits produced
an error ledger of confirmed errors in the deployment's outputs.
Study~C scores blind re-audits of the pre-audit project against that
ledger; Study~A tests the specification and the checklist on four
analysis tasks with known answers; Study~B tests the checklist on held-out
analyses from other projects. The corrected MaxToki-217M analyses re-run
most of the affected pipelines with corrected code and inputs.}
\label{fig:overview}
\end{figure}

\subsection*{The framework}

\subsubsection*{Pipeline specifications}
Each method is described in one markdown file. The template asks for
ten sections, in this order: overview; source (citation and pinned code
commit); inputs; outputs; dependencies; methodology (numbered steps with
equations and hyperparameters); code references (pointers into the pinned
reference code); a parameter table with defaults and ranges; validation
criteria; and known pitfalls. Only one of the eight deployed
specifications (the ageing-probe one) has all ten. The other seven have
no separate code-references section. Their code pointers are mostly inside
the methodology section, and they end with a quick-start section on
applying the method to a new model (see \emph{What deterministic checks
catch} in Results). The manifold specification names no pinned commit,
because its source paper links no code. The specifications for the eight
pipelines were written on 15 April 2026, at the start of the deployment,
by Claude Opus~4.7 under the author's direction, from the source papers
and their code; they are 48,000--96,000 characters long. During the deployment the ten-section
structure was an instruction to the agent; no code read or checked the
specifications. A specification therefore worked as a long, structured
prompt that the agent was asked to follow. The executing agent chose, for
example, the number of null-model draws and the number of cells where
compute was limited, and recorded some but not all of these choices.

\subsubsection*{Per-model project folders}
Model-specific code (loading the checkpoint, mapping genes to tokens,
encoding cells as model input) lives in a project folder separate from the
specifications, so that a specification can be pointed at a new model
without being rewritten. Each pipeline run writes its scripts, logs,
outputs and an agent-written run summary into its own folder.

\subsubsection*{The ten-pattern audit checklist}
The checklist (Table~\ref{tab:checklist}) was distilled from six
peer-review comment sets on the author's earlier interpretability papers.
For each pattern it gives the objection as a reviewer would state it, a
diagnostic signature to search for in run outputs and code, and a
prescription. The deployed version is a full audit specification.
Besides the ten patterns, it has the other template sections and a
step-by-step audit workflow: list the runs, check each pattern in each
run, compare runs, grade severity, propose actions and write a report.
For each pattern it also lists examples from the MaxToki project and cites
the review documents that raised it. Studies~A and B, and one arm of
Study~C, used a shorter \emph{generic} version, about 40\% of the length
of the deployed file. It keeps the ten pattern entries, without the
project-specific examples and the citations, and with a few
project-specific phrases made general. It adds a short opening
instruction and leaves out the workflow and the other sections. Another
arm of Study~C received the deployed file. Both versions are in the code
repository.

% TABLE: checklist (tab:checklist) is inserted here in the assembled file.

\subsubsection*{What was checked by code, by agent judgment, or by a person}
Table~\ref{tab:checks} lists each requirement the framework imposes and
how it was checked during the deployment. Apart from pass/fail flags that
two pipelines computed from their own numbers, every check was an agent
judgment; the author's role was to start runs, answer questions, and ask
whether issues had been addressed. To make the distinction testable we
wrote a deterministic checker (see \emph{Deterministic checks} below) and
report what it can and cannot detect.

% TABLE: checks (tab:checks) is inserted here.

\subsection*{The MaxToki-217M deployment}

\subsubsection*{Model and data}
MaxToki-217M~\cite{maxtoki2026} is a Llama-architecture autoregressive
transformer with 11 transformer blocks, hidden size 1,232, 8 attention
heads per block and a vocabulary of 20,275 tokens, 20,271 of them genes (Hugging Face
\texttt{theodoris-lab/MaxToki}). The downloaded weight files match
revision \texttt{21aa7b7f} byte for byte. That was the \texttt{main}
branch on the download day. The same bytes are also at earlier and later
revisions, so the hashes confirm the weights but not the exact revision.
The model exposes 12 hidden states, called layers 0--11 below. Layer~0
is the token embedding. For $l$ from 0 to 10, layer~$l$ is the input of
transformer block~$l$ (blocks counted from 0). Layer~11 is the output of
the last block after the final normalisation. Layer-8 attention is the
attention inside block~8. One attention run also used MaxToki-1B. Every pipeline gave
the model one cell per sequence (about 1,000--2,900 tokens per cell); no
multi-cell trajectory inputs were used. Cells came from the Replogle
essential-gene Perturb-seq screens in K562 and
RPE1~\cite{replogle2022perturbseq} (the K562 cells from the filtered public
copy in the Hugging Face dataset \texttt{arcinstitute/State-Replogle-Filtered},
and the RPE1 cells from the scPerturb copy~\cite{peidli2024scperturb}),
the Adamson CRISPR-interference
screen~\cite{adamson2016perturbseq} (in the processed form distributed with
GEARS~\cite{roohani2024gears}), Tabula
Sapiens~\cite{tabulasapiens2022} and, for the external-lung panel of the
topological hypothesis screen only, 1,500 cells from the Krasnow Lab Human
Lung Cell Atlas~\cite{travaglini2020lung} (Smart-seq2 data; CELLxGENE dataset
version \texttt{c88e0403-da93-40f4-99b5-f5fdeb81a82c}). Reference sets were
TRRUST~v2~\cite{han2018trrust}, ChIP-seq target sets from
DoRothEA~\cite{garcia-alonso2019dorothea}, the STRING protein-interaction
network~\cite{szklarczyk2023string}, and GO
biological-process~\cite{ashburner2000go}, KEGG~\cite{kanehisa2000kegg}
and Reactome~\cite{gillespie2022reactome} gene sets. In the
deployment, DoRothEA was used only in the audit's repairs of the
transcription-factor test; the corrected analyses use it for the same
test. STRING was used by the spectral-geometry, topology and SAE-atlas
runs. The gene sets were used mainly to annotate SAE features. The GO,
KEGG and Reactome sets are the Enrichr libraries
GO\_Biological\_Process\_2023, KEGG\_2021\_Human and Reactome\_2022, as
named in the script of an earlier project that built them. The
spectral-geometry and SAE-atlas runs also read GO terms, including
cellular-component terms, from a gene-to-GO table that appears to be the
one distributed with GEARS~\cite{roohani2024gears}. No file records which
release of DoRothEA or STRING, or which version of that gene-to-GO table,
was used.

\subsubsection*{Pipelines, compute and supervision}
Table~\ref{tab:pipelines} lists the eight pipelines and what was run.
All model computation of the deployment and of the corrected analyses ran
on one Apple-silicon MacBook Pro laptop with
32~GB unified memory. It used the PyTorch Metal (MPS) backend in 32-bit
floating point; some analysis ran on the same machine's CPU. The run
files do not name the chip. The machine that holds the files today has
an Apple M2 Pro and is probably the same one, but no file proves it.
From the first to the last run file the deployment spanned 21.9 days (15 April to 7 May 2026), with
files written on 11 of those days. Log files cover 47.8~hours of
computation when overlapping runs are counted once. The agent was Claude
Opus~4.7, run through Claude Code (Anthropic). The run files do not
record the model, and the session transcripts were not kept; the model is
known from the author, and a usage note that the agent wrote on 8 May
2026 also names it. The one automated hypothesis loop, in the
spectral-geometry pipeline, called Claude Sonnet (version not recorded).
The author typed 73
prompts during the deployment. Supervision time was not logged.

% TABLE: pipelines (tab:pipelines) is inserted here.

\subsubsection*{The in-session audit}
On the last day of the deployment the agent wrote the audit specification
and applied the checklist to all eight pipelines in one working session.
Earlier that day the same session had extended the topological hypothesis
screen, so for that pipeline it graded its own same-day work. All other
runs, including the first run of that screen, had been made in other
sessions. The agent wrote an 8~$\times$~10
matrix of free-text verdicts (pipeline $\times$ pattern) and a list of
proposed repairs, and then carried out two rounds of repair, each started
by a prompt from the author. From the audit request to the last repair
took 3.7~hours. The audit sweep read mainly the run summaries, plus
selected output tables and scripts, and re-ran no model. Three repair
steps did run MaxToki-217M again: a re-extraction of the lung cells with a
new sampling seed, a targeted re-run of the transcription-factor test for
GATA1, and a small cell-level benchmark.

\subsection*{Independent verification and the error ledger}
After the deployment, each claim in the run summaries and each number
they report was checked in four ways. (i)~A deterministic number tracer
searched each run's stored outputs for every number in its summary,
together with a chance rate for such matches. (ii)~Key statistics were
re-derived from stored per-record outputs with new code (for example,
direction-of-change predictions were rescored against constant-sign
baselines on the identical records; null models were re-sampled and
checked for actual randomisation). (iii)~Intervention code was read
against the transformer library source and unit-tested: an edit of zero
must leave the output unchanged, an edit at layer~$l$ must leave earlier
layers unchanged, and combined edits must differ from single edits.
(iv)~The model-input encoding of every run was compared with the
encoding the model expects. Claude Opus~5.5 agents ran checks (i)--(iv)
in working sessions that the author started and directed. Some of these
checks followed up points raised by a human reader from outside the
project (not the author). So an error credited to independent verification was found
by an agent doing a directed check, not by an agent working on its own.
Further errors came from the blind re-audits
of Study~C: every finding that matched no ledger entry was checked against
the full project by a separate agent, which confirmed or rejected it with
code. This agent's confirming code was not kept; the ledger rows give
the recomputed values. Four errors were first pointed out by the human reader.
Independent verification then confirmed each of them from the stored
files and code.

Each confirmed error is one row of the ledger (S1~Table), with its
location, the deterministic evidence that confirms it, its type (code
bug, wrong statistic, missing baseline or null, wrong description,
unsupported claim, reproducibility), the checklist pattern it falls
under (or ``outside the checklist''), whether it was present before the
in-session audit (from file timestamps), who or what first found it, and
whether the in-session audit report caught it. The full ledger in the code
release also records who or what confirmed each error. Severity was graded as \emph{critical} (a headline result or claim cannot
stand), \emph{major} (it stands only with material qualification) or
\emph{minor}. The ledger covers the deployment's code, logs, outputs, run
summaries and audit reports.

\subsection*{Subject agents and context control}
All agents in the pre-registered parts of the three controlled studies
(executors, reviewers, repairers and graders) were Claude Opus~5.5
(\texttt{claude-opus-5-5}). In the exploratory arm of Study~A (described
below), the executors, reviewers and repairers were Claude Haiku~4.5 and
the graders were Opus~5.5. All agents were run as Claude Code workflow
subagents of the built-in read-only type with the model set explicitly.
We used this agent type because it starts
without project instruction files and without the author's memory notes,
both of which contain descriptions of earlier errors; we checked this with
a probe agent before the studies. The agents can read files and run
shell commands. They were asked to run all analyses as inline Python
(\texttt{numpy}, \texttt{pandas}, \texttt{scipy}, \texttt{scikit-learn};
CPU only, at most 9~minutes per command) without writing files, and to
read only inside their task folder. Deliverables were returned as
structured output and written to disk by the harness. Study~A task folders and Study~B item folders were placed
at a neutral path under opaque names, so that folder names revealed
neither the study, the arm nor the status of an item. After each study we
scanned every tool call of the subject agents in its pre-registered part:
341 executors, reviewers, repairers and re-auditors in all (200 in
Study~A, 132 in Study~B and 9 in Study~C). Graders were meant to read the
keys and were not scanned. We looked for any access to the project
repository, the answer keys, the protocol files other than the checklist
given to a reviewer, the author's notes, this manuscript, or the folder
of another task or item. No subject agent accessed any of them. The
subject agents of the exploratory Haiku arm were checked in a separate
scan (see below and S2~Appendix). We also checked whether any agent
read a file that another agent had written. Twenty-five agents used a
scratch folder that the session shares between agents. Three of them
wrote temporary files there. The other 22 only ran a command to make sure
the folder existed. Six of these 22 also listed the folder and so saw the
names of other agents' files; none opened such a file.

In Study~A task T1 (whether layer-8 attention encodes TRRUST
regulator--target pairs in RPE1 cells), a Claude
Haiku~4.5 executor from the first run of the exploratory arm left a
script in the task folder. That run was set aside (see \emph{Study A}
below). All ten Opus~5.5 repair agents of T1 in the contract arm (the
agents that revised a deliverable after its review) opened this script. The contract arm is the Study~A arm in which the
executor also received the deployed specification. All ten of their
deliverables reached the correct verdict. Leaving out the T1
contract-arm runs does not change any Study~A result. For example, a
checklist review followed by repair still raised the share of pitfalls
handled by 2.9 percentage points over no review (95\% bootstrap interval
over the remaining 35 executor runs, 0.7 to 5.1). Graders of these
deliverables also opened that script.
Every task and item folder was afterwards checked to be identical to its
source.
Reasoning effort was the same session default for all agents.

The design, outcomes and analyses were written down before the studies
were run, and the answer keys were frozen by hash before the studies that
used them (S2~Appendix). We call this protocol pre-registered in that
sense: it was frozen by file hash and logged before the studies, but it
was not lodged in a public registry. The protocol has three amendments, all given in
S2~Appendix. Amendment~1 was logged on 1 October 2026, after Study~C had run
and before Studies~A and B started. It replaced one Study~A task, split
Study~A into two batches with the same protocol, and excluded one Study~B
pair that failed its own verification. Amendment~2 added the exploratory
Haiku arm described below, after the results of Study~A task T2 were
known. Amendment~3, logged before that arm was re-run, set aside the
arm's first run and gave every agent its own folder.

\subsection*{Study C: blind re-audit of the deployment}
We used file timestamps to rebuild a snapshot of the project as it stood
before the in-session audit started. The snapshot held all run code, logs
and configuration files, the specifications and the run summaries. It
also held most outputs smaller than 5~MB, and nine outputs of 5--20~MB
chosen because they hold the per-gene-pair or per-cell-group data behind
headline numbers. Every other output was replaced by a listing of its
name, size and shape. Where the audit had edited a summary in place, we
removed its dated audit additions. The audit files, all manuscript folders and
every later file were excluded. The reference set was the 68 ledger
errors that had been logged when the study started and were present in
this snapshot (6 critical, 25 major, 37 minor). When the set was frozen,
60 of them had been confirmed. The other 8 (3 critical, 3 major, 2 minor)
were confirmed later: 6 by new analyses, 1 by reading the code and 1 by
a literature check. All 68 are in S1~Table. Three arms received
the same project description, output format and budget instruction
(``Work as a careful expert would in about one working day. You may use
up to roughly 150 tool calls''): the checklist as deployed (the full
audit specification, with its project-specific examples and its step-by-step audit workflow, which
the prompt asked the agent to follow), the generic checklist (the ten
pattern entries only), or an instruction to review the project critically
before publication. Each arm was run three times.

Two graders, given the ledger and the audit's findings, judged for every
ledger error whether the audit detected it (\emph{yes}: a finding
identifies the same problem clearly enough that acting on it would lead
to the fix; \emph{partial}: it points at the same result or code with a
related but wrong or incomplete concern; \emph{no}), and matched every
finding to ledger rows. A third agent resolved disagreements. Findings
that matched no ledger row were checked by a fourth agent against the
snapshot and classified as real error, real but trivial, false alarm, or
undecidable; real errors were then verified against the full repository
and added to the ledger. The in-session audit was scored on the same 68
errors from its written sweep report, made before the repair rounds. The
Study~C graders did not score it. Its score is the ledger's record of
whether that report caught each error (yes, partly or no). The
independent verification made this record when it built the ledger,
before Study~C ran. The record was not blinded, and no second grader
checked it. ``Partly'' means that the report flagged
the area but not the actual error.

\subsection*{Study A: controlled execution study}
Four analysis tasks were built from the deployment's data
(Table~\ref{tab:tasks}). Each task asks a question about MaxToki-217M
that the deployment had answered, gives the data needed to answer it,
and asks for a report with a verdict, estimates with uncertainty, and a
complete analysis script. Each task has known pitfalls into which the
deployment fell (for example, a 50\% baseline for a direction-of-change
score whose classes are unbalanced). The pitfalls are not mentioned in
the task brief.

% TABLE: tasks (tab:tasks) is inserted here.

In the \emph{paper} arm the agent received the brief, the data and the
source paper of the method. In the \emph{contract} arm it also received
the deployed specification and one sentence asking it to follow it (the
arm names are those of the pre-registration, S2~Appendix). Both arms
therefore had the same task information and the same access to the
source method. Each executor output was then (a)~left as is,
(b)~reviewed by a fresh agent instructed to review it critically, or
(c)~reviewed by a fresh agent instructed to apply the generic checklist;
in (b) and (c) a fresh agent then revised the deliverable given the
review points, with the same repair instruction in both arms. Review
prompts had the same budget sentence (``Spend about the effort of a
careful two-hour expert review; you may use up to roughly 60 tool
calls.''). Each task--arm cell had five
independent executor runs, giving 40 executor runs and 120 final
deliverables.

Each answer key states the correct verdict and any conditions under which
a hedged verdict also counts as correct, key numbers with tolerances,
six pitfalls, each with what correct and incorrect handling looks like,
and tempting false statements. Each key was written by one agent. Its key numbers were computed by a
script that reads only the task's data. A second agent re-derived the key
numbers with its own code, and the keys were then frozen by hash. Two graders scored each deliverable against the key;
a third grader decided when the first two disagreed on the primary
outcome. Before grading, words that could reveal the arm (such as
``specification'', ``review'' and ``repair'') were replaced by a
placeholder in the deliverable. Graders also read the task brief, which in
the contract arm points to the specification file, so they could tell the
executor arm. Only the review arm was hidden from them.

The primary outcome was a correct verdict. Secondary outcomes were the
share of pitfalls handled (yes~=~1, partly~=~0.5, no~=~0), the number of
materially false statements, key numbers within tolerance, the share of
review points that a grader judged invalid, verdict changes caused by
repair, and whether the returned analysis script reproduced the reported
estimates. For this last outcome, the harness ran each script on the
task data with the same Python interpreter that the subjects used, with a
20-minute limit. Scripts that stopped while other scripts shared the
machine were run again, alone or two at a time, with a 60-minute limit;
none of these needed more than 10~minutes. A reported estimate counted as
reproduced when the script printed a number of the same name that agreed
with it within 0.005 or within 1\% of its value, whichever was larger.

There were four primary contrasts. The first compared the contract and
paper arms without review. It used a Fisher exact test, the risk
difference with a Newcombe interval~\cite{newcombe1998interval}, and a
Mantel--Haenszel test stratified by task. The other three compared review
arms on the same executor outputs: checklist versus generic, checklist
versus none, and generic versus none. These used exact McNemar
tests~\cite{mcnemar1947note}. We applied Holm
correction~\cite{holm1979simple} across all four contrasts.

Because every Opus~5.5 deliverable reached the correct verdict on the
first task to finish (T2), we added an exploratory arm after that task's
results were known, while the other tasks were still running
(Amendment~2): the same design with Claude Haiku~4.5 as executor, reviewer
and repairer, and the same Opus~5.5 graders and keys. A first run of this
arm was set aside and not analysed. Its Haiku agents did not keep to the
instruction not to write files: they wrote scripts into the shared
temporary folder and into shared task folders, where other agents could
read them, and some did. The arm was then run again with one folder per agent, each holding
links to the task files (Amendment~3, S2~Appendix). Only the re-run is
reported. S2~Appendix describes the checks of which files its agents
read. Five reviewers opened the folder of the executor whose deliverable
they were reviewing, because the deliverable named it. Every read of a
file in the shared temporary folder was of a file that the reading agent
had written itself. The arm is reported separately and does not enter the
pre-registered contrasts.

\subsection*{Study B: held-out audit benchmark}
Study~B used analyses that played no part in building the checklist.
They came from three other projects by the author. One studied
Cell2Sentence (C2S-Scale-Gemma-2-2B)~\cite{rizvi2025c2sscale} cell
embeddings. One studied the ESM-2 (650M)~\cite{lin2023esm2} protein
language model. One studied frozen scGPT, Geneformer V2-316M and
STATE-SE~\cite{adduri2025state} cell embeddings, together with
gene-regulatory-network and perturbation benchmarks, including the
Perturb-seq data of Norman et al.~\cite{norman2019exploring} as
distributed by scPerturb~\cite{peidli2024scperturb}. One pair used the
fetal human gut atlas of Elmentaite et al.~\cite{elmentaite2020gut}, and
three used proteins and residue annotations from
UniProtKB/Swiss-Prot~\cite{uniprot2025}. Four of the 11 pairs
in the study use no model at all (two label-building steps and two
benchmarks). Each analysis was rebuilt as a self-contained package (data,
code, outputs and a results summary written as the original analyst would
have written it) in a \emph{flawed} version with exactly one documented
error and a \emph{clean} version that differs only in that error. Of
twelve planned pairs, nine had a natural error that had been found and
fixed after the fact. The other three had a planted overclaim of causal
or mechanistic meaning in the summary (the checklist's pattern P9, for
which the pool had no natural example). Every package was checked by
code: the flawed version reproduces its wrong result, the clean version
the correct one, and comments or names that hinted at the error were
removed. One natural pair failed this check and was dropped, leaving 22
items (11 flawed, 11 clean; 6 errors inside the checklist's patterns and
5 outside, such as a circular statistic that breaks on evenly spread
angles or a row-index join that matched the wrong cells).

Each item was reviewed three times with the generic checklist and three
times with a generic critical review, with the same budget sentence in
both arms (``Spend about the effort of a careful one-hour expert review;
you may use up to roughly 50 tool calls.''). Two graders with the item's key judged
whether a review detected the documented error under a pre-written
detection rule, and classified every finding as the documented error,
another real problem, a false alarm, a minor or style point, or disputed;
a third grader resolved disagreements, and a verifier agent resolved
disputed findings by running code on the package. There were three
outcomes. The first was sensitivity: the share of reviews of flawed items
that found the documented error. The second was the share of reviews of
clean items with at least one false alarm. A finding counted as a false
alarm in one of two ways. In the first, the graders judged that it
claimed a problem that does not exist, checked against the item's key and
the data, and that it said a headline result or claim was wrong. In the
second, a grader marked the finding as disputed and the verifier agent,
after running code on the package, ruled it a false alarm. These findings
counted at any severity, whether or not they concerned a headline result.
The third was the number of other real problems found per review.

After the results were known, we also sorted the false alarms about a
main result that the graders had marked into four groups: correct but
changing no stated conclusion; correct and already stated in the
analysis's own summary or README; factually wrong; or a real problem that
the key had missed. This step was not pre-registered. Each claim was
labelled by one agent, with no second rater. False alarms that the
verifier agent ruled after a grader dispute were not sorted.

\subsection*{Corrected MaxToki-217M analyses}
The verification found that the intervention code of two pipelines (SAE
circuit tracing, and exhaustive mapping and steering) edited the wrong
tensor and that every run except the RPE1 attention run encoded the model
input incorrectly (Results). We therefore re-ran the main affected
analyses (those not re-run are listed in Results) with three changes.
First, the model input follows MaxToki's rank-value encoding. Each gene's count is divided by that gene's median in the
Geneformer gene-median table (gc104M) that the MaxToki tokenizer uses,
and genes are ranked from high to low. The K562 and Adamson files store
only $\log(1+\text{CP10k})$ values, where CP10k means counts scaled so
that each cell sums to 10,000. For these files we undid the log with
$\exp(x)-1$, divided by the smallest non-zero value in the cell and
rounded. This gives whole-number counts up to one factor per cell, so the
gene ranking is the same as from the raw counts. For the RPE1, Tabula
Sapiens and lung files we used the stored raw counts. Second, the
intervention hooks edit the input of transformer block~$l$, which is the
tensor that the layer-$l$ sparse autoencoder reads. They add each edit to
the live activations, so edits at several layers add up instead of
overwriting each other. Third, we retrained all twelve sparse
autoencoders (one per layer state) on correctly encoded cells. Every
model input was checked against the expected token order before each
forward pass. Each re-run that edits the model had sanity checks (zero
edit gives zero effect; edits change only later layers; combined edits
differ from single ones; a feature that never fires changes nothing). The
key numbers of each corrected analysis were computed a second way with
separate code. Other agents also re-derived them with their own code:
three for the RPE1 attention analysis and one for the cross-model
comparison (S3~Appendix). Methods
specific to each analysis are given with its result and in S3~Appendix.

\subsection*{Deterministic checks}
The checker has three parts: a specification linter (presence and order
of the ten sections, a resolvable pinned commit, code references in the
required form, a machine-readable parameter table), a run-manifest check
(parameters used versus the specification's defaults and ranges, gate
flags recomputed from stored values) and the number tracer described
above. It reads files only and runs no model. Its code and outputs are in
the repository.

\subsection*{Statistics}
Proportions are reported with Wilson 95\% intervals~\cite{wilson1927probable}
unless stated otherwise. Bootstrap intervals are percentile intervals.
The one exception is the feature-triplet analysis, which uses basic
bootstrap intervals over cells, as stated with its result. The resampling
unit is the item in Study~B, the reference error in Study~C, the executor run in Study~A, the
deployed pipeline run (eight in all) for the number tracer, and the knockdown
(silenced gene), transcription factor, gene, donor or cell in the
MaxToki-217M analyses. The spectral-geometry analysis uses jackknife
intervals instead, over blocks of cells or groups of genes. In Study~A
the main intervals resample executor runs pooled over all four tasks. The
pre-registered plan resamples runs within each task. Results report
intervals of that kind as well, for the main study and for the Haiku~4.5
arm. Study~A bootstrap intervals use 4,000 draws; those of Studies~B and
C use 2,000. In the MaxToki-217M analyses, multiple testing within a
family of tests was controlled with the Benjamini--Hochberg (BH)
procedure~\cite{benjamini1995fdr}. A BH-adjusted $p$-value is called a
$q$ value, and a test passes when $q < 0.05$. Study~A used Holm
correction instead, as described above. Empirical $p$-values use the
$(b+1)/(m+1)$ form~\cite{phipson2010permp}. Degree-preserving network
nulls used the Curveball algorithm~\cite{strona2014curveball} and were
checked to change the network in every draw. The corrected analyses ran
in Python 3.12.9 with PyTorch 2.11.0, transformers 5.5.4, NumPy 1.26.4,
SciPy 1.17.1, pandas 2.3.3 and scikit-learn 1.8.0, as recorded in their
run configuration files; the deployment's runs did not record package
versions. All code is at
\url{https://github.com/Biodyn-AI/maxtoki-mechinterp}.

\subsection*{Declaration of generative AI and AI-assisted technologies}
Generative AI was the object of this study and also a tool in it. The
deployment (specifications, pipeline execution, run summaries and the
in-session audit) was carried out by Claude Opus~4.7 through Claude Code,
except for one automated loop that called Claude Sonnet, as described
above. The controlled studies used Claude Opus~5.5 as subjects
and graders. Opus~5.5 agents also wrote the Study~A answer keys (a
second agent re-derived their key numbers before the keys were frozen),
classified review findings that matched no ledger entry, checked disputed
Study~B findings by running code, and did the post-hoc labelling of false
alarms, as described above. The protocol, the Study~A tasks, the Study~B
packages, the deterministic checker and the supporting files were drafted
by Opus~5.5 agents under the author's direction and checked by separate
agents and by code. The exploratory arm of Study~A used Claude Haiku~4.5 as
executor, reviewer and repairer, with Opus~5.5 graders, as described
above. Claude Opus~5.5 also carried out the independent verification, the
corrected analyses (the key numbers of each were computed a second way
with separate code; other agents also re-derived them, three for the RPE1
attention analysis and one for the cross-model comparison), the statistical
analysis code and the figure code, and
drafted this manuscript under the author's direction. The author reviewed
and edited all content and takes full responsibility for it. AI tools were
not used to create or alter any image of primary data; all figures are
plotted by code from the stored outputs.
